\appendix
\section{Matching radial moments and the heat exterior}
\label{app:moments}

The leading swirl $E(X,\eta)$ has, by \ref{P3}, the far-field
\[
E(X,\eta)=c_\infty X^{-1/2-h}\bigl(1+O(X^{-1})\bigr),
\qquad X\to\infty.
\]
In physical variables, on $z=0$, this is
\[
u^{(0)}_\theta(r,0,t)
=q^{-A} c_\infty \Bigl(\frac{r^2}{2q}\Bigr)^{-1/2-h}\bigl(1+O(q/r^2)\bigr)
=c_\infty' r^{-1-2h} q^{h}\bigl(1+O(q/r^2)\bigr).
\]
As $q\downarrow 0$ at fixed $r>0$ this tends to $0$. The heat exterior $\Gamma(r,t)$ is chosen so that its radial moments
\[
M_n(t):=\int_0^\infty \Gamma(r,t)\, r^{n}\, r\,dr
\]
match the corresponding moments of $u^{(0)}_\theta$ through the annulus, for $n=0,1,\dots,N$ with $N$ large. Because $\Gamma$ solves the radial heat equation of Proposition~\ref{prop:heat}, each $M_n$ solves a closed ODE. Prescribing the jet $\{M_n(t)\}_{n\le N}$ at each $t$ is a finite-dimensional interpolation problem, solved by a linear combination of Gaussians in $r^2$ with time-dependent coefficients. Sending $N\to\infty$ in a Borel-summable way produces a remainder decaying faster than every power of $q$ in the matching region, which is absorbed into $\mathcal E_{\infty}$.

The same construction works with a slight $z$-dependence in a buffer layer, at the cost of additional heat kernels in the axial variable; because axial diffusion is smaller by $q^{2h}$, this buffer is cheap.

\section{Analytic profiles near the axis}
\label{app:axis}

\subsection{Even-odd structure}

A Cartesian-smooth axisymmetric field with swirl admits a power series in $r^2$ near the axis:
\[
u_z=\sum_{m\ge 0} a_m(z,t) r^{2m},
\qquad
\frac{u_\theta}{r}=\sum_{m\ge 0} b_m(z,t) r^{2m},
\qquad
\frac{u_r}{r}=\sum_{m\ge 0} c_m(z,t) r^{2m}.
\]
In similarity variables this is a power series in $X$. Substituting into the leading operator $\mathcal P$ of Lemma~\ref{lem:rescaled-op} yields a recursive system for the coefficient functions of $\eta$. The first coefficients are constrained by \ref{P4}: $U(0,0)\neq 0$ means $a_0\neq 0$ at $z=0$, $\tau$ small.

\subsection{Shooting}

On a rectangle $[0,X_*]\times[-\eta_c,\eta_c]$ we consider the Banach space $\mathcal B$ of pairs $(E,U)$ whose axis jets satisfy the Cartesian smoothness conditions, with norm
\[
\|(E,U)\|_{\mathcal B}
=\sup_{0\le k\le k_0}\Bigl(
\|X^{-1/2}E\|_{C^k}+\|U\|_{C^k}
\Bigr).
\]
The map that sends a pair to the solution of the linearised leading balance with that pair as Dirichlet data at $X=X_*$ (where we impose the beginning of the far-field \ref{P3}) is a contraction for $h$ small and $k_0$ large, because the operator $\mathcal P$ is a regular perturbation of a radial {S}tokes operator with a lower-order swirl coupling of size $h$. We omit the expansion of the contraction constants in this draft and record the conclusion as Proposition~\ref{prop:profiles}. The far-field $E\sim c_\infty X^{-1/2-h}$ is the unique decaying solution of the constant-coefficient system obtained by linearising $\mathcal P$ at $X=\infty$ with $U=V_0=0$; matching at $X_*$ fixes $c_\infty>0$.

\subsection{Axial bias}

The reflection $z\mapsto -z$, $\eta\mapsto -\eta$, $U\mapsto -U$ is a symmetry of the leading equations. We break it by imposing $U(0,0)=\beta$ for a small $\beta\neq 0$. For $\beta=0$ the plane $\eta=0$ is simultaneously a nodal plane of $U$ and a plane of weak rotational amplification. For $\beta\neq 0$ these loci separate by $O(\beta)$, which is the content of \ref{P4} and \ref{P6} near $z=0$. Continuity in $\beta$ preserves \ref{P1}--\ref{P3} and \ref{P5}.

\section{The admissible stress cone}
\label{app:cone}

Let $e_r,e_\theta,e_z$ be the cylindrical frame. A pulse polarisation $\Psi$ that is divergence-free to WKB order, with local wavevector in the $(r,z)$-plane, has vanishing $\theta$-mean. Its quadratic mean is a symmetric tensor of the form
\[
S(\Psi)
=\overline{\Psi\otimes\Psi}
=\alpha\, e_r\otimes e_r
+\beta\, (e_r\otimes e_\theta+e_\theta\otimes e_r)
+\gamma\, (e_r\otimes e_z+e_z\otimes e_r)
+\cdots,
\]
where the coefficients $(\beta,\gamma)$ are the radial fluxes of azimuthal and axial momentum. As the polarisation and the wavevector vary in the open set permitted by \ref{P6}, the pair $(\beta,\gamma)$ covers an open cone $\mathcal K_0\subset\R^2$. We write $\mathcal K$ for the cone of symmetric tensors whose $(r\theta,rz)$ flux pair lies in $\mathcal K_0$ and whose remaining components are determined by the WKB polarisation (they are not independent).

\begin{lemma}[Two families span the cone]
\label{lem:cone}
There exist two polarisations $\Psi_+,\Psi_-$ such that the corresponding flux pairs $(\beta_+,\gamma_+)$ and $(\beta_-,\gamma_-)$ are linearly independent. Consequently every interior point of $\mathcal K_0$ is a unique positive combination $\iota_+(\beta_+,\gamma_+)+\iota_-(\beta_-,\gamma_-)$ with $\iota_\pm>0$.
\end{lemma}

\begin{proof}
This is two-by-two linear algebra, once two independent flux ratios have been exhibited. Family $+1$ is chosen with a dominant $r\theta$ flux (a ``swirl pulse''); family $-1$ with a dominant $rz$ flux (an ``axial pulse''). Independence is arranged by \ref{P6}, which provides both rotational and axial shear as energy sources, and therefore both polarisations as unstable modes of the geometric-optics problem.
\end{proof}

The {B}ogovskii lift used in Proposition~\ref{prop:Tann} produces a stress in a neighbourhood of $u_B\otimes u_B$ restricted to the annulus. For $q$ small this neighbourhood lies in the interior of $\mathcal K$, because $u_B\otimes u_B$ itself has a definite $(r\theta,rz)$ signature on the support of $\chi_{\ann}$, by the signs of $E$ and of the axial bias.

\section{Index of exponents}
\label{app:exponents}

For the reader's convenience we collect the exponents at viscosity $1$, with $\tau=1-t$ and $0<h<1/100$.

\begin{center}
\begin{tabular}{@{}lll@{}}
\toprule
Quantity & Symbol & Order \\
\midrule
Radial core width & $\ell_r$ & $\tau^{1/2}$ \\
Axial core height & $\ell_z$ & $\tau^{1/2-h}$ \\
Concentration scale & $q$ & $\tau$ on the core \\
Similarity radius & $X$ & $r^2/(2q)$ \\
Azimuthal / axial speed & $|u_\theta|,\,|u_z|$ & $\tau^{-1/2-h}$ \\
Radial speed & $|u_r|$ & $\tau^{-1/2}$ \\
Pressure & $|p|$ & $\tau^{-1-2h}$ \\
Core energy & $\|u\|_{L^2(\mathcal C_\tau)}^2$ & $\tau^{1/2-3h}$ \\
Angular {R}eynolds number & $\Ree_\theta$ & $\tau^{-h}$ \\
Radial {R}eynolds number & $\Ree_r$ & $O(1)$ \\
Axial / radial diffusion & $\ell_r^2/\ell_z^2$ & $\tau^{2h}$ \\
Annular residual & $|\Rmom{u_B}{p_B}|$ & $\tau^{-3/2-h}$ \\
Annular stress & $|T_{\ann}|$ & $\tau^{-1-2h}$ \\
\bottomrule
\end{tabular}
\end{center}

\noindent
The blowup ring used throughout is $z=0$, $r=\sqrt{2X_*\tau}$, on which $u_\theta=E(X_*,0)\tau^{-1/2-h}$.
